\begin{remark}
\label{remark-topologies}
In terms of topologies
Lemmas \ref{lemma-dominate-etale-neighbourhood-finite-flat} and
\ref{lemma-dominate-etale-affine-finite-flat} mean the following.
Let $S$ be any scheme. Let $\{f_i : U_i \to S\}$ be an \'etale covering
of $S$. There exists a Zariski open covering $S = \bigcup V_j$,
for each $j$ a finite locally free, surjective morphism
$W_j \to V_j$, and for each $j$ a Zariski open covering
$\{W_{j, k} \to W_j\}$ such that the family
$\{W_{j, k} \to S\}$ refines the given \'etale covering
$\{f_i : U_i \to S\}$. What does this mean in practice?
Well, for example, suppose we have a descent problem which we
know how to solve for Zariski coverings and for fppf coverings
of the form $\{\pi : T \to S\}$ with $\pi$ finite locally free
and surjective. Then this descent problem has an affirmative
answer for \'etale coverings as well. This trick was used by
Gabber in his proof that $\text{Br}(X) = \text{Br}'(X)$
for an affine scheme $X$, see \cite{Hoobler}.
\end{remark}